\documentclass[addpoints]{exam}
% \documentclass[addpoints,12pt]{exam}

\usepackage[slovene]{babel}
\usepackage{amsfonts}
\usepackage{amssymb}
\usepackage{amsmath}
\usepackage{float}
\usepackage{graphicx}
\usepackage{enumitem}
\usepackage{color}
\usepackage[gen]{eurosym}
\usepackage{newunicodechar}
\newunicodechar{€}{\euro}

%Komentar naslednje vrstice glede na verzijo testa -- rešitve ali prostor za odgovore:
% \printanswers

% \linespread{2}


\pointsinrightmargin
\marginbonuspointname{}


\renewcommand{\solutiontitle}{\noindent\textbf{Rešitve in točkovanje:}\par\noindent}

\colorgrids
\definecolor{GridColor}{gray}{0.7}
\setlength{\gridsize}{7mm}

\extrawidth{5mm}
\setlength{\rightpointsmargin}{1.5cm}

\begin{document}

\runningfooter{}{}{}

\begin{center}
    \fbox{\fbox{\parbox{15cm}{\centering
    Popravljanje in pridobivanje ocen -- konferenca \\ 1. letnik -- test B \\ 15. junij 2026 \\ (Racionalna števila; Koordinatni sistem v ravnini; Funkcija; Premica)}}}
\end{center}

\vspace{0.1in}
\makebox[\textwidth]{Ime in priimek, oddelek:\enspace\hrulefill}


\begin{center}
    \hqword{Naloga}
    \hpword{Možne točke}
    \chbpword{Dodatne točke}
    \hsword{Dosežene točke}
    \htword{Skupaj}  
    % \combinedpointtable[h][questions]
    \gradetable[h][questions]

\end{center}

\vspace{0.02in}


\begin{questions}

    %%%%% racionalna števila
    
    % 1. vprašanje

    \question[4]
    Izračunajte natančno vrednost izraza.
    % \begin{parts}
    %     \part[4]
        $$\dfrac{35\cdot \frac{1}{7}+\frac{7}{8}}{2\frac{1}{3}+\frac{7}{4}}-4\frac{8}{9}:\frac{11}{3}$$

        \begin{solution}[9cm]
            \begin{align*} 
                \dfrac{\frac{2}{3}+1\frac{5}{6}\cdot 3}{\frac{2}{7}:\frac{4}{3}-1\frac{1}{14}+2}&=\dfrac{\frac{2}{3}+\frac{11}{6}\cdot 3}{\frac{2}{7}\cdot\frac{3}{4}-\frac{15}{14}+\frac{28}{14}}= \\
                                                                                                &=\dfrac{\frac{2}{3}+\frac{11}{2}}{\frac{1}{7}\cdot\frac{3}{2}+\frac{13}{14}}= \\
                                                                                                &=\dfrac{\frac{4}{6}+\frac{33}{6}}{\frac{3}{14}+\frac{13}{14}}= \\
                                                                                                &=\dfrac{\frac{37}{6}}{\frac{16}{14}}= \\
                                                                                                &=\dfrac{37\cdot 14}{6\cdot 16}= \\
                                                                                                &=\dfrac{259}{48} 
            \end{align*}
            Za pravilno seštevanje/odštevanje ulomkov $1$ točka; pravilno množenje ulomkov $1$ točka; pravilno razrešen dvojni ulomek $1$ točka, okrajšan rezultat $1$ točka.
        \end{solution}




        \question[4]
        Kolikšno koncentracijo raztopine dobimo, če $2~kg$ $60~\%$ raztopine razredčimo z $2~kg$ $35~\%$ raztopine?

        \begin{solution}[5cm]
            \begin{align*}
                0.60 \cdot 2~kg + 0.40\cdot 3~kg &= x \cdot 5~kg \\
                                 1.2~kg+1.2~kg &= x\cdot 5~kg \\
                                            x &= \dfrac{2.4~kg}{5~kg} \\
                                            x &= 0.48
            \end{align*}
            Pravilna nastavitev enačbe $1$ točka, pravilno reševanje $1$ točka, pravilen rezultat $1$ točka, zapis pravilnega odgovora $1$ točka.
        \end{solution}





    \newpage

        % 3. vprašanje

        \question
        Poenostavite izraze.
        \begin{parts}
            

            \part[5]
            $\dfrac{1}{x+2}-\dfrac{3}{x-2}-4\left(4-x^2\right)^{-1}$
    
            \begin{solution}[8cm]
                \begin{align*}
                    \left(\dfrac{6-x}{x^2+6x}-\dfrac{x}{36-x^2}\right)\left(\dfrac{2x-6}{x^2+6x}\right)^{-1}+\dfrac{x}{6-x} &= \left(\dfrac{6-x}{x(x+6)}-\dfrac{x}{(6-x)(6+x)}\right)\dfrac{x(x+6)}{2(x-3)}+\dfrac{x}{6-x}= \\
                                                                                                                            &= \left(\dfrac{(6-x)^2}{x(x+6)(6-x)}-\dfrac{x^2}{x(6-x)(6+x)}\right)\dfrac{x(x+6)}{2(x-3)}+\dfrac{x}{6-x}= \\
                                                                                                                            &= \dfrac{36-12x+x^2-x^2}{x(x+6)(6-x)}\dfrac{x(x+6)}{2(x-3)}+\dfrac{x}{6-x}= \\
                                                                                                                            &= \dfrac{12(3-x)}{x(x+6)(6-x)}\dfrac{x(x+6)}{2(x-3)}+\dfrac{x}{6-x}= \\
                                                                                                                            &= \dfrac{-6}{6-x}+\dfrac{x}{6-x}= \\
                                                                                                                            &= \dfrac{-6+x}{6-x}= \\
                                                                                                                            &= -1
                \end{align*}
                Za pravilno faktorizacijo izrazov $1$ točka, pravilno razčlenjevanje izrazov $1$ točka, pravilno seštevanje ulomkov $1$ točka, pravilna obratna vrednost ulomka $1$ točka, pravilno množenje ulomkov $1$ točka, pravilno krajšanje ulomkov $1$ točka, okrajšan rezultat $1$ točka.
            \end{solution}
    
    
    
    
            \part[5]
            $\dfrac{\left(9x^3z^2\right)^{-4}}{-3^4x^{-2}y^3} : \left(\dfrac{3^3x^2z^{-2}}{-y^{-3}}\right)^{-6}$    
            
            \begin{solution}[7cm]
                \begin{align*}
                    \left(\dfrac{-3^4x^{-2}y^3}{x^3z^2}\right)^{-4} : \left(\dfrac{y^{-3}}{3^5x^2z^{-2}}\right)^3 &= \dfrac{3^{-16}x^{8}y^{-12}}{x^{-12}z^{-8}} : \dfrac{y^{-9}}{3^{15}x^6z^{-6}}= \\
                                                                                                                    &= \dfrac{3^{-16}x^{8}y^{-12}}{x^{-12}z^{-8}} \cdot \dfrac{3^{15}x^6z^{-6}}{y^{-9}}= \\
                                                                                                                    &= 3^{-16+15}x^{8-(-12)+6}y^{-12-(-9)}z^{-(-8)+(-6)}= \\
                                                                                                                    &=\frac{x^{26}z^2}{3y^3}
                \end{align*}
                Za pravilno potenciranje z negativnim eksponentom $1$ točka, pravilno potenciranje s pozitivnim eksponentom $1$ točka, pravilno deljenje $1$ točka, pravilno krajšanje $1$ točka, končen okrajšan rezultat $1$ točka.
            \end{solution}
    
    
        \end{parts}
   
   
   
   



        \question
        Ceno puloverja so zvišali za $20~\%$, a ker ni šel v prodajo, so ga nato pocenili za $30~\%$.
        Po drugi pocenitvi ga je kupil Matej in zanj plačal $30.24~\euro$.

        \begin{parts}
            \part[1]
            Koliko odstotkov prvotne cene puloverja je plačal Matej?
            \begin{solution}[0.5cm]
                
            \end{solution}

            % \part[2]
            % Izračunajte začetno ceno puloverja?
            % \begin{solution}[5cm]
                
            % \end{solution}

            \part[2]
            Določite ceno puloverja neposredno pred drugim znižanjem?

        \end{parts}

        \begin{solution}[2cm]
            \begin{align*}
                0.70\cdot \left(1.10\cdot x\right) &=115.5~\euro \\
                                     0.77 \cdot x &= 115.5~\euro \\
                                                x &= \dfrac{115.5}{0.77}~\euro \\
                                                x &= 150~\euro
            \end{align*}
            Pravilen zapis in reševanje enačbe $1$ točka, pravilen rezultat $1$ točka, zapis pravilnega odgovora $1$ točka.
        \end{solution}

        

%%%% PKS, funkcija, premica

    % 1. vprašanje

    \question[4]
    Narišite dano množico točk v ravnini.
        $$[-1,2)\times \{-1,1,2\}$$

            \begin{solution}[9cm]
                ????
            \end{solution}

    
    
    % \newpage

    % 2. vprašanje
    \question[4]
    Izračunajte ploščino in orientacijo trikotnika $\triangle ABC$ z oglišči v točkah $A(4,-3)$, $B(5,4)$ in $C(-4,1)$.
    
    
    \begin{solution}[6cm]
        ???
    \end{solution}

        


\newpage

    % 5. vprašanje
    \question
    Dana je funkcija $$f(x)=\begin{cases}
    -2x-2, & x\leq 0 \\ -1, & 0<x<1 \\ 0.5x+1, & x\geq 2
    \end{cases}.$$ 
    
    \begin{parts}
        \part[6]
        Narišite dano funkcijo.

            \begin{solution}[11cm]
                ???
            \end{solution}

        \part[1]
        Določite definicijsko območje funkcije $f$.
        
            \begin{solution}[1cm]
                ???
            \end{solution}

        \part[3]
        Izračunajte ničle in začetno vrednost funkcije $f$.

            \begin{solution}[6cm]
                ???
            \end{solution}

        % \part[2]
        % Določite in utemeljite, ali je funkcija surjektivna. 
        
        %     \begin{solution}[3cm]
        %         ???
        %     \end{solution}

        \part[1]
        Določite interval, kjer je funkcija $f$ konstantna.
        
            \begin{solution}[0cm]
                ???
            \end{solution}

    \end{parts}
    
    






\end{questions}



\end{document}